%!TEX root = ../csuthesis_main.tex

{~}
\vspace{-9pt}
\section{攻读学位期间主要的研究成果} % 无章节编号

\ifblindreview
% \noindent
% （盲审隐去作者相关具体信息）
\fi
\vspace{11pt}
\subsection*{一、学术论文}

\ifblindreview

% \noindent
% 第一作者：JCR 1 区 x 篇，会议 x 篇 \\{}
% 第二作者：JCR 1 区 x 篇，3 区 x 篇，4 区 x 篇，EI x 篇 

% \noindent
% 投稿状态： 
% IEEE Transactions on Image Processing 1篇（under review）\\{}
% IEEE Transactions on Circuits and Systems for Video Technology 1 篇（Accept with Minor Revision）

% 学位办老师要求用如下这种几乎算是单盲的格式，我也木有办法……
%\noindent
\begin{enumerate}[label={[\arabic*]},itemindent=2em,wide]
	\item 
	
\end{enumerate}

\else
% 标准版本
\begin{enumerate}[label={[\arabic*]},itemindent=2em,wide]
	\item \textbf{Wanru Dai}, Wenxuan Wang, Yuhu Lu, Peng Du, Jinjing Shi et al. Quantum audio neural networks with time-series encoding for audio classification[J]. Quantum Information Processing, 2026, 25: 142. {（SCI收录，JCR1区，影响因子：2.2，本人一作）} 
	
\end{enumerate}
\fi

\vspace{22pt}
\subsection*{二、发明专利}
\ifblindreview
\begin{enumerate}[label={[\arabic*]},itemindent=2em,wide]
	\item 
\end{enumerate}

\else
\begin{enumerate}[label={[\arabic*]},itemindent=2em,wide]
	\item 石金晶，戴婉如，基于音频编码的量子卷积神经网络模型. 2026.01.06，中国，2025112891325. (实审阶段)
	
	
	
	
\end{enumerate}
\fi

\ifblindreview
\vspace{22pt}
\subsection*{三、主持和参与的科研项目}
%\noindent
\begin{enumerate}[label={[\arabic*]},itemindent=2em,wide]
	\item[(1)] 
\end{enumerate}
\else

% {~}
\vspace{22pt}
\subsection*{三、主持和参与的科研项目}
%\noindent
\begin{enumerate}[label={[\arabic*]},itemindent=2em,wide]
	\item[(1)] 国家自然科学基金项目：参数化量子线路学习建模的实用方法和技术，项目编号：62272483，时间：2023.09-2026.05，参与.
	\item[(2)] 湖南省自然科学基金杰出青年基金项目：基于玻色采样的量子机器学习理论与方法，项目编号：2023JJ10078，时间：2023.01-2025.12，参与.
\end{enumerate}
\fi

\ifblindreview
\vspace{22pt}
\subsection*{四、个人获奖情况}
%\noindent
\begin{enumerate}[label={[\arabic*]},itemindent=2em,wide]
	\item[(1)] 
\end{enumerate}
\else
% {~}
\vspace{22pt}
\subsection*{四、个人获奖情况}
%\noindent
\begin{enumerate}[label={[\arabic*]},itemindent=2em,wide]
	\item[(1)] 校级，中南大学一等学业奖学金，2023.
	\item[(2)] 校级，中南大学二等学业奖学金，2024.
	\item[(3)] 院级，中南大学电子信息学院研究生学术年会优秀论文，2025.
\end{enumerate}
\fi
\ifblindreview
\else

\newpage
\section{{致~~~~谢}} % 无章节编号
% 重新设置正文行间距，因为前置部分设置时候行间距被改过
\renewcommand*{\baselinestretch}{1.0} % 几倍行间距
\setlength{\baselineskip}{20pt} % 基准行间距
{
	
	时光荏苒，行文至此，二十余载求学路暂告段落。回首在中南大学的数年时光，感慨万千，感激满怀。
	
	首先，谨以最诚挚的敬意感谢我的导师石金晶教授。从论文的选题立意、框架搭建，到研究过程中的疑难解惑，直至最终的修改定稿，无不倾注了您大量的心血。您渊博的学识、严谨的治学态度、敏锐的学术洞察力以及宽厚谦和的为人风范，令我深深敬仰，是我终生学习的榜样。师恩如海，铭记于心。
	
	感谢同组的施鹤远老师、胡超老师和冷健老师，感谢你们在课程学习与论文开题、答辩过程中给予的宝贵指导与无私帮助。同时，衷心感谢实验室、课题组的全体同门。特别感谢王雯萱师姐、陆玉虎师兄、赵仁鑫师兄、陈添师姐和肖子萌师兄在我初入课题时的悉心引领与慷慨分享；感谢吕绚丽和张程阳同学，与你们并肩作战、互相扶持的日子，是科研道路上最温暖的陪伴；也感谢杜鹏、李世成、王硕蕾、杨永祥、陈睿、谌思，你们的陪伴与支持，给我的科研生活带来了许多欢乐与帮助。感谢这个集体赋予我的所有指导、灵感与情谊。
	
	深深感谢我的父母及家人。你们是我最坚实的后盾，是我前行路上不竭的动力源泉。你们的无私付出、默默支持与无限包容，让我能心无旁骛地完成学业。养育之恩，无以为报，唯有不断努力，成为你们的骄傲。
	
	感谢我的室友段唐慧、李佳欣、张钰霖，你们的陪伴、倾听与鼓励，为我驱散了无数迷茫与压力，让这段旅程充满色彩。
	
	最后，向百忙之中参与本论文评阅和答辩的各位专家、老师致以最衷心的感谢！由于本人学识有限，文中难免存在疏漏与不足，恳请各位师长批评指正。
	
	山水一程，三生有幸。新的旅程即将开始，我自当心怀感恩，砥砺前行。
}
\vspace{-9pt}
\vspace{1.5cm}

\hfill 2026年3月8日

\hfill 戴婉如于中南大学

\fi